\section{Related Literature}
\label{sec:literature}

% TODO for Zotero: confirm and attach PDFs for morris_1989, schwarz_lane_1998,
% nyc_charter_revision_commission_2025, been_madar_mcdonnell_2014,
% hankinson_2018, schleicher_2013, hills_schleicher_2011, and
% fang_stewart_tyndall_2023.

This project starts from the large literature showing that land-use regulation limits housing supply and raises housing costs. \citet{glaeser_why_2005} argue that regulation helps explain why Manhattan prices sit far above construction costs, and \citet{glaeser_economic_2018} frame this gap as an implicit regulatory tax. \citet{paciorek_supply_2013} shows that supply constraints matter not only for price levels, but also for how quickly markets respond to demand shocks. \citet{turner_land_2014} estimate the welfare costs of land-use regulation, while \citet{molloy_housing_2022} connect housing supply limits to affordability and household location. The common point is that housing supply is not just a technological margin. It is also a political and administrative margin.

The second strand asks why local governments restrict housing even when new supply has broad benefits. 
\citet{fischel_homevoter_2009} gives the canonical homeowner mechanism: owner-occupiers have large, undiversified stakes in local housing values and therefore have strong incentives to oppose changes they see as risky. 
\citet{ortalo-magne_political_2014} formalize how homeownership and anti-growth politics can reinforce each other. 
\citet{einstein_neighborhood_2019} show that the people who participate in local land-use meetings are not representative of the broader public, 
and \citet{hankinson_when_2018} shows that support for housing can change sharply with spatial scale. 
These papers help motivate the treatment in this project, which is baseline homeownership in a community district relative to the borough when the court case was decided.
Relative homeownership shares are a proxy for the political concentration of residents with the strongest incentives to resist nearby large buildings.

A newer set of papers studies how the scale of land-use authority changes housing outcomes. \citet{mast_warding_2024} finds that moving from at-large to ward elections reduces housing permits by shifting representation toward smaller local constituencies. 
\citet{khan_decentralized_2021} studies Chicago's aldermanic privilege and shows that local control can make zoning more restrictive near ward boundaries. 
\citet{fang_homeowner_2023} find that Toronto councillors representing more homeowners are more likely to oppose housing, especially large projects in their own wards. 
These papers are closest to this project because they treat political geography as part of the housing supply technology.

The New York City setting adds a different piece of variation. 
The institutional change was not a town changing electoral rules or a ward boundary moving. 
It was a court-induced charter reform that replaced the Board of Estimate with a district-based City Council at the final stage of land-use review. 
Legal scholarship on New York land use has long argued that seriatim, district-specific review can lead city officials to ignore the aggregate housing costs of local restrictions \citep{hills_jr_balancing_2011,schleicher_city_2013}. \citet{been_urban_2014} also show that New York City rezonings often reflect homeowner and neighborhood politics rather than a simple growth-machine story. 
This paper connects those institutional accounts to a measurable prediction: after final authority shifts to a district-based council operating under member deference, large-building construction should fall most in community districts where homeowner opposition was likely to be strongest.
